
\(\xi\) la classe d'un tel faisceau. On a alors :
\[
 \ch(\lambda_{-1}(\xi))=\ch(1-\xi)=1-\ch(\xi).
\]
Comme \(\widetilde c(\xi)=(1,1+c^1(\xi))\), on obtient, d'après la définition de \(\ch\) :
\[
 \ch(\lambda_{-1}(\xi))=1-\exp(c^1(\xi)).
\]

Par ailleurs, \(\xi\) est d'augmentation \(1\); la classe de Chern qui intervient dans la formule 4.4 relative à \(\xi\) est donc \(c^1\). Compte tenu de la définition de \(\Todd\), et de 3.5, la relation à démontrer s'écrit :
\[
 1-\exp(c^1(\xi))
 =
 -c^1(\xi)\cdot\left(\frac{-c^1(\xi)}{1-\exp(c^1(\xi))}\right)^{-1},
\]
et est donc vérifiée.

\textbf{Remarque 4.4.1.} Il est possible aussi de déduire 4.4 d'un énoncé purement algébrique correspondant, cf. \(\RRR\) I 1.3.

\textbf{4.5.} Nous allons montrer maintenant que pour montrer le théorème de Riemann--Roch dans le cas envisagé dans ce paragraphe, on peut se ramener au cas où \(Y\) est un diviseur de \(X\), et où l'élément auquel s'applique la formule est image inverse d'un élément de \(K^{\bullet}(X)\).

D'après 4.2, la formule à démontrer :
\[
 \ch(i_*(x))=
 i_{\mathrm{gr}}\bigl(\ch(x)\cdot\Todd(\check T_i)\bigr)
\]
est équivalente à la formule :
\[
 f^{\mathrm{gr}}\bigl(\ch(i_*(x))\bigr)
 =
 f^{\mathrm{gr}}\circ i_{\mathrm{gr}}\bigl(\ch(x)\cdot\Todd(\check T_i)\bigr).
\]
Compte tenu de la fonctorialité de \(\ch\), et de 4.3, elle peut encore s'écrire :
\[
 \ch\bigl(f^*(i_*(x))\bigr)
 =
 j_{\mathrm{gr}}\bigl(g^{\mathrm{gr}}\bigl(\ch(x)\cdot\Todd(\check T_i)\bigr)\cdot c^{d-1}(\check F)\bigr),
\]
soit, d'après VII 3.4 :
\[
 \ch\bigl(j_*(g^*(x)\cdot\lambda_{-1}(F))\bigr)
 =
 j_{\mathrm{gr}}\bigl(\ch(g^*(x))\cdot\Todd(g^*(\check T_i))\cdot c^{d-1}(\check F)\bigr).
\]

